\section{Introduction}
\label{sec:intro}

Fine-tuning a safety-aligned language model on a narrow, seemingly innocuous task can
break its alignment in broad and surprising ways. \citet{betley2025em} showed that
training a model to write \emph{insecure code} --- with no mention of malice in the data ---
induces \emph{emergent misalignment}: the model later gives harmful advice and voices
anti-human views on completely unrelated prompts. This is a striking safety failure. But it
is also a striking \emph{partial} failure: the misaligned behavior appears on only about
20\% of evaluation answers, and many fine-tuning runs produce no broad drift at all. The
inconsistency is not noise to be averaged away --- it is the scientific signal. It implies
that models hold \emph{latent robustness regions} where adversarial fine-tuning simply fails
to move behavior.

{\bf Why does this matter?} Knowing \emph{which} behaviors resist drift, and \emph{why},
is central to both AI safety (which capabilities are protected) and adversarial research
(which are exploitable). A natural and appealing hypothesis is that robustness is
\emph{graded by semantic distance}: a prompt far from the fine-tuning domain --- a math word
problem, a science question --- should be insulated from a code-focused intervention,
forming a low-drift ``basin'' on orthogonal tasks. This intuition is implicit in the
``safety basin'' of \citet{peng2024landscape} and in the orthogonal-direction analyses of
\citet{yang2025asft}.

{\bf What is missing?} Prior work characterizes \emph{aggregate} safety at the
\emph{endpoint} of fine-tuning, or maps basins in \emph{weight} space. No prior work tests
the \emph{graded correlation between behavioral robustness and semantic distance}, measured
at the level of \emph{per-prompt output trajectories across training checkpoints}.
\citet{springer2026geometry} even argue the opposite of the basin intuition: orthogonality
to the fine-tuning objective is structurally unstable and collapses under gradient descent.
Whether behavioral robustness is graded, and whether it is a stable attractor, is an open,
testable question.

{\bf Our approach.} We introduce a checkpoint-level \emph{trajectory-divergence} diagnostic
(\method). We LoRA-fine-tune a real instruct model, \smol, on the canonical \insecure-code
dataset of \citet{betley2025em}, checkpointing during training. At each checkpoint we greedily
decode a \emph{graded probe set} of 15 prompts spread across 5 semantic-distance tiers, from
in-domain code to orthogonal math and science. We measure how far each probe's output drifts
from its own pre-fine-tuning baseline using a composite of embedding cosine distance and
character edit distance, and we test whether drift falls with semantic distance. Crucially,
we repeat the entire pipeline on a matched benign \secure-code control, isolating
adversarial \emph{intent} from generic fine-tuning drift.

{\bf What we find.} The simple graded-robustness picture fails, and in an informative way.
Drift does \emph{not} decrease with semantic distance (Spearman $\rho=-0.08$, $p=0.78$,
$n=15$). Orthogonal-capability prompts drift the \emph{most} (mean 0.55), while a conceptual
code-safety Q\&A pocket drifts the \emph{least} (0.045) --- the reverse of the basin
prediction. Most decisively, the adversarial and benign fine-tunes produce nearly identical
drift (global means 0.356 vs.\ 0.352; every per-tier intent difference $\le 0.061$). At this
scale, trajectory divergence reflects generic fine-tuning style, not adversarial intent. The
one reproducible low-drift pocket we find is content-type-specific (factual/technical
knowledge preserved nearly verbatim), not task distance.

{\bf Contributions.}
\begin{itemize}[leftmargin=*,itemsep=0pt,topsep=0pt]
    \item We propose \method, a checkpoint-level trajectory-divergence diagnostic that
    measures per-prompt behavioral drift from a pre-fine-tuning baseline, paired with a
    graded probe set ordered by semantic distance from the fine-tuning domain.
    \item We conduct a controlled test of the graded-robustness hypothesis on a real instruct
    model, and refute it at this scale: drift is non-monotonic in distance, with orthogonal
    tasks drifting most and a factual-knowledge pocket drifting least.
    \item We show, by differencing against a matched benign-fine-tuning control, that raw
    trajectory drift is dominated by generic fine-tuning style rather than adversarial intent
    --- a concrete methodological correction for using such metrics as robustness signals.
    \item We release our pipeline, per-probe trajectories, and analysis to support scaling
    the diagnostic to larger models where emergent misalignment is known to appear.
\end{itemize}

The rest of the paper is organized as follows. \Secref{sec:related} positions our work
against the fine-tuning-safety, alignment-geometry, and training-dynamics literatures.
\Secref{sec:method} describes \method, the graded probe set, and the divergence metric.
\Secref{sec:results} reports the core hypothesis tests, and \secref{sec:discussion}
interprets them, states limitations, and draws implications.
